\exercisegroup

\begin{enumerate}
\def\labelenumi{\arabic{enumi})}
\item
  Let G be a graph and let \(Q = (S, F, o, t)\) its underlying quiver.
\end{enumerate}

\begin{enumerate}
\def\labelenumi{\alph{enumi})}
\item
  Show that one has \(d_{+}(s) = d_{-}(s)\) for all \(s \in S\). This integer is then called the degree of G at s and denoted \(d(s)\).
\item
  Assume that G is finite. For there to exist a closed path \((a_{0}, f_{1}, \ldots, f_{n}, a_{n})\) in G such that, for every edge \(\varphi\) of G, there exists a unique integer n such that \(\varphi = \{f_{n}, \overline{f_{n}}\}\), it is necessary and sufficient that G be connected and that \(d(s)\) be even for all \(s \in \text{Som}(G)\).
\end{enumerate}

\begin{enumerate}
\def\labelenumi{\arabic{enumi})}
\setcounter{enumi}{1}
\tightlist
\item
  Let \(\mathrm{G} = (\mathrm{S}, \mathrm{F}, o, t)\) be a quiver. We say that G is a bipartite quiver if there exists a partition \(P = \{S_{1}, S_{2}\}\) of S into subsets such that, for every arrow \(f \in F\), the origin \(o(f)\) and the term \(t(f)\) are in different parts of the partition P.
\end{enumerate}

\begin{enumerate}
\def\labelenumi{\alph{enumi})}
\item
  Assume that G is connected. Prove that there exists at most one partition P as above.
\item
  Prove that a graph G is bipartite if and only if every cycle in G has even length.
\end{enumerate}

\begin{enumerate}
\def\labelenumi{\arabic{enumi})}
\setcounter{enumi}{2}
\item
  Let G be a connected, locally finite graph whose vertex set is infinite. Prove that there exists a composable sequence \((f_{n})_{n\in\mathbf{N}}\) of arrows of G whose origins are pairwise distinct (König's theorem).
\item
  Let G be a finite graph and let \(Q = (S, F, o, t)\) be its underlying quiver. We endow the space \(C^{S}\) with the unique Hermitian space structure for which its canonical basis is orthonormal. Let L: \(C^{S} \to C^{S}\) be the map given by \(\mathrm{L}(u)(x) = \sum_{o(f)=x} (u(t(f)) - u(x))\) for \(x \in S\) and \(u \in C^{S}\).
\end{enumerate}

\begin{enumerate}
\def\labelenumi{\alph{enumi})}
\item
  Prove that the endomorphism L of \(C^{S}\) is positive (EVT, V, p. 45, def. 6).
\item
  Prove that the kernel of L consists of the functions \(u \in C^{S}\) which are constant on each connected component of G.
\item
  Let \(F_{1}\) be an orientation of G. Let \(E = (e_{s,f})\) be the matrix of type \((\mathrm{S}, \mathrm{F}_{1})\) given by
\end{enumerate}

\[
 e _ {s, f} = \left\{ \begin{array}{l l} 1 & \text {if } o (f) = s \text {and } t (f) \neq s; \\ - 1 & \text {if } o (f) \neq s \text {and } t (f) = s; \\ 0 & \text {otherwise. } \end{array} \right.
\]

Prove that the matrix \(\Lambda\) of L in the canonical basis of \(\mathbf{C}^{\mathrm{S}}\) is given by \(\Lambda = E\cdot {}^{t}E\) .

\(d)\) Assume that G is connected and that its set of vertices is nonempty; let \(x\) be a vertex of G and let \(\mathrm{S}' = \mathrm{S} - \{x\}\). Let T be a subset of F of cardinality \(\mathrm{Card}(\mathrm{S}')\). Prove that T is the set of arrows of a maximal oriented subtree of G (equipped with the orientation \(\mathrm{F}_1\)) if and only if the rank of the matrix \(E_{\mathrm{T}}'\) of type \((\mathrm{S}',\mathrm{T})\) deduced from \(E\) is equal to \(\mathrm{Card}(\mathrm{S}')\). From this, deduce that the determinant of the matrix \(A_{\mathrm{S}',\mathrm{S}'}\) is equal to the number of maximal oriented subtrees of G. (Use exercise 6 of A, III, §8, p. 192.)

\begin{enumerate}
\def\labelenumi{\alph{enumi})}
\setcounter{enumi}{4}
\tightlist
\item
  Let W be the orthogonal complement of \(\mathrm{Ker}(\mathrm{L})\). Prove that W is a subspace of \(\mathbf{C}^{\mathrm{S}}\) that is stable under W and that \(\det (\mathrm{L}|_{\mathrm{W}})\) is equal to the number of maximal forests of the graph G (Kirchhoff's theorem).
\end{enumerate}

\begin{enumerate}
\def\labelenumi{\arabic{enumi})}
\setcounter{enumi}{4}
\tightlist
\item
  \begin{enumerate}
  \def\labelenumii{\alph{enumii})}
  \tightlist
  \item
    Let n be a natural integer \(\geqslant 1\) and let G be a graph whose set of vertices is the set \(\mathbf{Z}/n\mathbf{Z}\) and such that two vertices i and j are joined by an arrow if and only if \(i-j=\pm1\pmod n\). Describe explicitly the set of maximal forests of G.
  \end{enumerate}
\end{enumerate}

\begin{enumerate}
\def\labelenumi{\alph{enumi})}
\setcounter{enumi}{1}
\item
  Let G be a graph and let S be the set of its vertices. Suppose that, for every pair \((x,y)\) of vertices of G, the set of arrows of G with origin x and terminus y has cardinality 1 if \(x \neq y\) and 0 otherwise ("complete graph"). Compute the cardinality of the set of maximal forests of G.
\item
  Let G be a graph and let S be the set of its vertices. Let \((\mathrm{S}_{1},\mathrm{S}_{2})\) be a partition of S; suppose that, for every pair \((x,y)\) of vertices of G, the set of arrows of G with origin x and terminus y has cardinality 1 if \((x,y)\in\mathrm{S}_{1}\times\mathrm{S}_{2}\) and 0 otherwise ("complete bipartite graph"). Compute the cardinality of the set of maximal forests of G.
\end{enumerate}

\begin{enumerate}
\def\labelenumi{\arabic{enumi})}
\setcounter{enumi}{5}
\tightlist
\item
  Let \(c\) be a natural integer.
\end{enumerate}

\begin{enumerate}
\def\labelenumi{\alph{enumi})}
\tightlist
\item
  There exists a unique family of maps \(\mathrm{R}_{c}\colon(\mathbf{N}^{*})^{c}\to\mathbf{N}^{*}\), for \(c\geqslant2\), satisfying the following properties:
\end{enumerate}

\begin{enumerate}
\def\labelenumi{(\roman{enumi})}
\item
  If there exists i such that \(n_{i}=1\), then \(\mathrm{R}_{c}(n_{1},\ldots,n_{c})=1\);
\item
  We have \(\mathrm{R}_{2}(m,n)=\mathrm{R}_{2}(m-1,n)+\mathrm{R}_{2}(m,n-1)\) if m and n are integers \(\geqslant2\);
\item
  If \(c \geqslant 3\), we have \(\mathrm{R}_c(n_1, \ldots, n_c) = \mathrm{R}_{c-1}(n_1, \ldots, n_{c-2}, \mathrm{R}_2(n_{c-1}, n_c))\) for all \((n_1, \ldots, n_c) \in \mathbf{N}^c\).
\end{enumerate}

\begin{enumerate}
\def\labelenumi{\alph{enumi})}
\setcounter{enumi}{1}
\item
  We have \(\mathrm{R}_2(m,n) = \binom{m + n - 2}{m - 1}\) for every pair \((m,n)\) of integers \(\geqslant 1\).
\item
  For every finite sequence \((n_{1},\ldots,n_{c})\) of natural integers \(\geqslant1\) and every natural integer \(R\geqslant1\), we say that the property \(\boldsymbol{R}(n_{1},\ldots,n_{c};\mathrm{R})\) is satisfied if, for every complete graph whose set of vertices S has cardinality \(\geqslant R\) and every partition \(P=(F_{1},\ldots,F_{c})\) of the set F of its arrows, there exists an integer \(i\in\{1,\ldots,c\}\) and a subset A of S of cardinality \(n_{i}\) such that every arrow of G connecting two elements of A belongs to \(F_{i}\).
\end{enumerate}

Prove the property \(\boldsymbol{R}(n_{1},\ldots,n_{c};\mathrm{R}_{c}(n_{1},\ldots,n_{c}))\) .

\begin{enumerate}
\def\labelenumi{\alph{enumi})}
\setcounter{enumi}{3}
\item
  We denote by \(\mathrm{R}(n_{1},\ldots,n_{c})\) the smallest integer \(R\geqslant1\) such that the property \(\boldsymbol{R}(n_{1},\ldots,n_{c};\mathrm{R})\) is satisfied (“Ramsey numbers”).
\item
  We have \(R(n,2)=n\) for every integer \(n\geqslant2\); we have \(R(3,3)=6\).
\item
  Prove that R(4,3) = 9.
\end{enumerate}

Compute R(5, 5); compute R(6, 6).

\begin{enumerate}
\def\labelenumi{\arabic{enumi})}
\setcounter{enumi}{6}
\tightlist
\item
  Let G be a group acting on the right on a set X, and let A be a subset of G. The Schreier graph \(\mathcal{G}_{\mathrm{G}}(\mathrm{X},\mathrm{A})\) is the quiver \((\mathrm{S},\mathrm{F},o,t)\) defined by:
\end{enumerate}

- The set of its vertices is S = X;

- The set of its arrows is F = X × A;

- The maps \(o\) and \(t\) are given by \(o(x, a) = x\) and \(t(x, a) = x \cdot a\) for all \((x, a) \in \mathrm{X} \times \mathrm{A}\) .

  When X = G, on which G acts by right multiplication, we denote this quiver by \(\mathcal{G}(G,A)\) and call it the Cayley quiver of G with respect to A.

We call the graph associated with this quiver a Schreier graph (resp. Cayley graph).

\begin{enumerate}
\def\labelenumi{\alph{enumi})}
\item
  Prove that the quiver \(\mathcal{G}_{\mathrm{G}}(\mathrm{X},\mathrm{A})\) is connected if and only if the subgroup H of G generated by A has at most one orbit on X.
\item
  Assume that A generates G. Show that the Cayley graph \(\mathcal{G}(\mathrm{G},\mathrm{A})\) is connected and that it is bipartite if and only if there exists a group homomorphism \(\varphi \colon \mathrm{G}\to \mathbf{Z} / 2\mathbf{Z}\) such that \(\varphi (a) = 1\) for all \(a\in \mathrm{A}\) .
\item
  Show that there exists a unique action of G on the Cayley graph \(\mathcal{G}(\mathrm{G},\mathrm{A})\) such that the action of G induced on the vertices of this graph is the action of G on itself by left multiplication.
\end{enumerate}

\begin{enumerate}
\def\labelenumi{\arabic{enumi})}
\setcounter{enumi}{7}
\tightlist
\item
  A cycle \((a_0, f_1, \ldots, f_n, a_n)\) in a graph G is said to be Hamiltonian if, for every vertex \(s\) of G, there exists a unique integer \(m \in \{1, \ldots, n\}\) such that \(a_m = s\). A graph is said to be Hamiltonian if it has a Hamiltonian cycle in G.
\end{enumerate}

\begin{enumerate}
\def\labelenumi{\alph{enumi})}
\item
  Let S be a finite set of cardinality \(\geqslant 3\) and let \(K_{S}\) be a complete graph with vertex set S. Prove that the graph \(K_{S}\) is Hamiltonian. More precisely, prove that for every orientation A of G, there exists a Hamiltonian cycle in \(K_{S}\) whose every arrow belongs to A.
\item
  We denote by \(G_{S}\) the set of subgraphs of \(K_{S}\) whose set of vertices is equal to S. Let \(\preccurlyeq\) be the coarsest order relation on the set \(G_{S}\) for which \(G \preccurlyeq G'\) if there exist elements u, v of S such that \(\deg(u) + \deg(v) \geqslant \text{Card}(S)\) and such that \(\text{Fl}(G')\) is the union of \(\text{Fl}(G)\) and the two arrows of \(K_{S}\) with endpoints u and v.
\end{enumerate}

  Prove that, for every \(G \in G_{S}\), the set of elements H of \(G_{S}\) such that \(G \preccurlyeq H\) has a unique maximal element; it is denoted \(c(G)\).

\begin{enumerate}
\def\labelenumi{\alph{enumi})}
\setcounter{enumi}{2}
\item
  Let G be a subgraph of \(K_{S}\) with vertex set S. Prove that the graph G is Hamiltonian if and only if the same is true of the graph \(c(\mathrm{G})\) .
\item
  Let G be a subgraph of \(K_{S}\) whose vertex set is S and whose degree at each vertex is \(\geqslant\) Card(S)/2. Prove that \(c(\mathrm{G})\) is equal to \(K_{S}\) . Deduce that G is a Hamiltonian graph (theorem of G. Dirac).
\end{enumerate}

\begin{enumerate}
\def\labelenumi{\arabic{enumi})}
\setcounter{enumi}{8}
\tightlist
\item
  Let (W, S) be a Coxeter system (LIE, IV, §1, p. 11, def. 3); suppose that the group W is finite.
\end{enumerate}

\begin{enumerate}
\def\labelenumi{\alph{enumi})}
\item
  Suppose that (W, S) is irreducible and has at most two vertices. Prove that the Cayley graph \(\mathcal{G}(\mathrm{W},\mathrm{S})\) is Hamiltonian.
\item
  Suppose that (W, S) is irreducible and has at least three vertices. Let \(s \in S\) be a terminal vertex (LIE IV, appendix) of the Coxeter graph associated with (W, S). Set \(S' = S - \{s\}\) and denote by \(W'\) the subgroup of W generated by \(S'\). Let \(\Gamma\) be the graph whose set of vertices is the set of left cosets modulo \(W'\) in W and such that two distinct vertices A and B are joined by an edge exactly if and only if \(\cap\) is not empty. Using a maximal subtree of the graph \(\Gamma\), prove that if the Cayley graph \(\mathcal{G}(W', S')\) is Hamiltonian, then so is the graph \(\mathcal{G}(W, S)\).
\item
  Prove that the Cayley graph \(\mathcal{G}(\mathrm{W},\mathrm{S})\) is Hamiltonian (Conway-Sloane-Wilks theorem).
\end{enumerate}

\begin{enumerate}
\def\labelenumi{\arabic{enumi})}
\setcounter{enumi}{9}
\tightlist
\item
  Let G be a connected graph whose set of vertices is S. For \(x, y \in S\), we set
\end{enumerate}

\(d_{\mathrm{G}}(x,y) = \inf \{\ell \in \mathbf{N}\colon \text{il existe un chemin de longueur} \ell \text{ reliant } x \text{ à } y\}.\)

\begin{enumerate}
\def\labelenumi{\alph{enumi})}
\item
  Prove that \(d_{\mathrm{G}}\) is a distance on the set S.
\item
  Let \(G_{1}\) and \(G_{2}\) be graphs and \(\varphi\colon G_{1}\to G_{2}\) a graph morphism. Show that \(d_{\mathrm{G}_{2}}(\varphi(x),\varphi(y))\leqslant d_{\mathrm{G}_{1}}(x,y)\) for every pair \((x,y)\) of vertices of \(G_{1}\) .
\end{enumerate}

\begin{enumerate}
\def\labelenumi{\arabic{enumi})}
\setcounter{enumi}{10}
\tightlist
\item
  Let G be a connected graph whose vertex set S is nonempty. We call the diameter of G the diameter \(\operatorname{diam}(G)\) of the metric space (S, \(d_{\mathrm{G}}\)).
\end{enumerate}

\begin{enumerate}
\def\labelenumi{\alph{enumi})}
\item
  Prove that \(\text{diam}(G) \leqslant \text{Card}(S)-1\). For which graphs is this inequality an equality?
\item
  Let v be an integer \(\geqslant 2\) such that every vertex of G has degree at most v. Prove that \(\text{Card}(S) \leqslant v^{\text{diam}(G)}\). For which graphs is this inequality an equality?
\item
  Let n be an integer such that \(n \geqslant 2\). Let \(T_{n}\) be the subset of \(S_{n}\) formed by the transpositions with support \(\{m, m + 1\}\), for \(m \in \{1, \ldots, n - 1\}\). Prove that the diameter of the Cayley graph \(\mathcal{G}(\mathfrak{S}_{n}, \mathrm{T}_{n})\) is equal to \(n(n - 1)/2\).
\item
  Let n be an integer such that \(n \geqslant 2\). Let \(\tau\) be the transposition with support \(\{1,2\}\) and \(\sigma\) be the cycle \((1,2,\ldots,n) \mapsto (2,3,\ldots,n,1)\) in the group \(S_n\). Let \(G_n\) be the Cayley graph \(\mathcal{G}(\mathfrak{S}_n,\{\sigma,\tau\})\). Prove that \(n(n-1)/6 \leqslant \text{diam}(G_n) \leqslant 3n(n-1)/2\). (Let T be the set of sequences \((x,y,z)\) of elements of \(\{1,\ldots,n\}\) such that \(x \textless{} y \textless{} z\). If \(t=(x,y,z)\in\text{T}\), we say that a permutation \(\lambda\in S_n\) preserves the cyclic order of t if one has \(\lambda(x)<\lambda(y)<\lambda(z)\), or \(\lambda(y)<\lambda(z)<\lambda(x)\), or \(\lambda(z)<\lambda(x)<\lambda(y)\). Observe that the permutation \(\sigma\) preserves the cyclic order of every element of T, and that the permutation \(\tau\) preserves the cyclic order of every element of T that is not of the form \((1,2,x)\). Deduce that the distance from the identity permutation to the permutation \((1,\ldots,n) \mapsto (n,n-1,\ldots,1)\) is at least equal to \(n(n-1)/6\).
\end{enumerate}

\begin{enumerate}
\def\labelenumi{\arabic{enumi})}
\setcounter{enumi}{11}
\tightlist
\item
  Let \(n\) be an integer \(\geqslant 2\) and let \(p\) be a prime number \(>n/2\). Let H be a subgroup of \(\mathfrak{S}_n\) that acts transitively on \(\{1, \ldots, n\}\), contains a transposition and a cycle of order \(p\). Let G be a graph such that
\end{enumerate}

- The set of its vertices is \(S = \{1, \ldots, n\}\) ;

- The set F of its arrows is the set of pairs \((i,j)\in\mathrm{S}\times\mathrm{S}\) such that the transposition \(\tau_{i,j}\) belongs to H;

- We have \(o((i,j)) = i\) , \(t((i,j)) = j\) and \(\overline{(i,j)} = (j,i)\) for all \((i,j) \in \mathrm{F}\) .

\begin{enumerate}
\def\labelenumi{\alph{enumi})}
\item
  Prove that every connected component of G is a complete graph.
\item
  Prove that if G is connected, then \(H = S_{n}\) .
\item
  Prove that there exists a group homomorphism from \(\mathfrak{S}_n\) to \(\operatorname{Aut}(\mathbf{G})\) such that, for every \(\sigma \in \mathfrak{S}_n\), the map \(\operatorname{Som}(\sigma)\) is equal to the permutation \(\sigma\). Then show that \(\mathfrak{S}_n\) acts transitively on the set of connected components of \(\mathbf{G}\), and that these are pairwise isomorphic.
\end{enumerate}

\(d)\) Let \(\sigma\in\mathrm{H}\) be a cycle of order \(p\). Prove that \(\sigma\) stabilizes each connected component of G. Conclude that \(\mathrm{H}=\mathfrak{S}_{n}\).

